\documentclass{beamer}
\usepackage{beamerthemeCambridgeUS}
\usepackage{color}  % for \textcolor{red}{blah}
\usepackage{graphicx}
\usepackage{subfigure}
\usepackage{multicol}
\usepackage{verbatim} % for \begin{comment}
\usepackage{lipsum}


\newcommand\Fontvi{\fontsize{6}{7.2}\selectfont}


\newenvironment{packed_enum}{
\begin{enumerate}
  \setlength{\itemsep}{1pt}
  \setlength{\parskip}{0pt}
  \setlength{\parsep}{0pt}
}{\end{enumerate}}

\newenvironment{packed_item}{
\begin{itemize}
  \setlength{\itemsep}{1pt}
  \setlength{\parskip}{0pt}
  \setlength{\parsep}{0pt}
}{\end{itemize}}

\usetheme{CambridgeUS}

\title{Discrete Mathematics}
\author[Geoffrey Buhl]{Geoffrey Buhl \\ \texttt{geoffrey.buhl@csuci.edu}}
\institute[CSUCI]{California State University Channel Islands}
\date{week 6 class 1}

\def \vec{{\rm vec}}
\newtheorem{conjecture}[theorem]{Conjecture}
\newtheorem{explanation}[theorem]{Explanation}


\begin{document}

%\frame{\titlepage}

%\section[Outline]{}
%\frame{\tableofcontents}

\section{Week 6 Class 2}

\frame
{
  \frametitle{Solving Recurrence Relations}

Today we develop a second method for solving recurrence relations, the generating function method.

\vfill
Class Session Goals:
\begin{packed_enum}
\item Present Generating Function method for solving recurrence relations
\item Practice Generating Function method for solving recurrence relations
\end{packed_enum}
\vfill
}

\frame
{
  \frametitle{Taylor Series}

The Taylor Series of a function $f(x)$ at $x=0$ is:\\

$$f(0)+f'(0)x+\frac{f''(0)x^2}{2}+\frac{f^{(3)}(0)x^3}{3!}+\frac{f^{(4)}(0)x^4}{4!}+ \cdots=\sum^{\infty}_{n=0}\frac{f^{(n)}(0)x^n}{n!}$$

where $f^{(n)}$ is the $n$th derivative of $f$ evaluated at zero.  If $f(x)$ is {\bf analytic} then $f(x)$ is equal to its Taylor series. That is, if $f$ is analytic

$$f(x)=\sum^{\infty}_{n=0}\frac{f^{(n)}(0)x^n}{n!}$$
}

\frame
{
  \frametitle{Important Taylor Expansions}


$$\frac{1}{1-x}=1+x+x^2+x^3+x^4+x^5+x^6+ \cdots=\sum_{n=0}^{\infty}x^n$$
$$\frac{1}{(1-x)^2}=1+2x+3x^2+4x^3+5x^4+6x^5+ \cdots=\sum_{n=0}^{\infty}(n+1)x^n$$
$$\frac{1}{1-\alpha x}=1+\alpha x+\alpha^2x^2+\alpha^3x^3+\alpha^4x^4+\alpha^5x^5+ \cdots=\sum_{n=0}^{\infty}\alpha^nx^n$$
 This last series is called the {\bf Geometric Series} and it will be particularly useful to us.

}

\frame
{
  \frametitle{Taylor Series Example and IDEA!}
You can obtain the following expansion by differentiating the Taylor Series for $\frac{1}{(1-x)^2}$ and multiplying by$\frac{x}{2}$:
$$\frac{x}{(1-x)^3}=1x+3x^2+6x^3+10x^4+15x^5+21x^6+ \cdots=\sum_{n=1}^{\infty}\frac{n(n+1)}{2}x^n$$
\pause AND!
$$\frac{x}{(1-x)^3}=T_1x+T_2x^2+T_3x^3+T_4x^4+T_5x^5+T_6x^6+ \cdots=\sum_{n=1}^{\infty}T_nx^n$$
where $T_n$ are the triangular numbers 1,3,6,10,15, $\ldots$.
}



\frame
{
  \frametitle{Generating Function Method Idea}

{\bf Step 1} Given a sequence $a_n$, create a power series $\sum^\infty_0a_nx^n$ \vfill
{\bf Step 2} Find a function $f(x)$ that satisfies $f(x)=\sum^\infty_0 a_n x^n$ \vfill
{\bf Step 3} Use $f(x)$ to obtain a closed form solution for $a_n$ using known Taylor series. \vfill

This is an overview.  There are many details that the next example will demonstrate.

}


\frame
{
  \frametitle{Generating Function Method Example}

Let $a_n$ count the number of is the number of ternary sequences of length $n$ where the numbers of 0's is even and the number of 1's is odd.  Then $a_n=2a_{n-1}+3a_{n-2}$ with $a_0=0$ and $a_1=1$  (A015518). \vfill

Calculation shortcuts
Partial fractions: $\frac{x}{1-2x-3x^2}=-\frac{1}{4}\cdot\frac{1}{1+x}+\frac{1}{4}\cdot\frac{1}{1-3x}$ \vfill
Geometric Series with $\alpha=-1$: $\frac{1}{1+x}=\sum_{n=0}^{\infty}(-1)^nx^n$\vfill
Geometric Series with $\alpha=3$: $\frac{1}{1-3x}=\sum_{n=0}^{\infty}3^nx^n$\vfill
\pause \underline{Closed form}\\$a_n=\frac{3^n - (-1)^n}{4}$

}

\frame
{
  \frametitle{Generating Function Method Example}

Let $a_n$ count the number of $n$-jeweled necklaces with a clasp where adjacent jewels must different made out of 3 colored jewels.  Then $a_n=a_{n-1}+2a_{n-2}$ with $a_1=0$ and $a_2=6$  (A092297). \vfill

\underline{Calculation shortcuts}\\
Partial fractions: $\frac{6x}{1-x-2x^2}=\frac{2}{1-2 x}+\frac{-2}{1+x}$ \vfill
Geometric Series with $\alpha=-1$: $\frac{1}{1+x}=\sum_{n=0}^{\infty}(-1)^nx^n$\vfill
Geometric Series with $\alpha=2$: $\frac{1}{1+x}=\sum_{n=0}^{\infty}2^nx^n$\vfill
\pause \underline{Closed form}\\$a_n=2^n - 2(-1)^{n-1}$

}

\frame
{
  \frametitle{Generating Function Method Activity}

In small groups, work on the ``Generating Function'' worksheet. The goals of this worksheet are to:
\vfill
\begin{packed_enum}
\item Practice finding the generating function of a sequence
\item Practice obtaining a closed form solution from the Taylor series of a generating function.
\end{packed_enum}
\pause
\vfill
\underline{Problem 1 answers}\\$f(x)=\sum_{n=0}^\infty\left( \frac{13}{5}4^n+\frac{12}{5}(-1)^n\right)x^{n} $\\$a_n=\frac{1}{5}\left(12 (-1)^n+13\cdot 4^n\right).$
\pause
\vfill
\underline{Problem 2 answers}\\ $f(x)=\sum_{n=0}^\infty\frac{1}{5}\left( 3^n-(-2)^n\right)x^{n+1}$\\ $b_n=\left(3^{n-1}-(-2)^{n-1}\right)$
}

\begin{comment}
\frame
{
  \frametitle{We need a recurrence (or two) }

\vfill
{\bf Problem 1}: Consider tilings $a_n$ of a $1\times n$ board with $1 \times 2$ red tiles, $1 \times 2$ white tiles, $1 \times 2$ blue tiles, $1 \times 1$ orange tiles, and $1 \times 1$ cyan tiles.  What is the recurrence relationship for $a_n$ and what are the initial conditions?
\vfill

{\bf Problem 2}: Consider Problem 1 but with $r$ different colors of $1 \times 2$ tiles and $s$ different colors of $1 \times 1$ blocks.  What is the recurrence relationship to describe the number of tilings $b_n$ of a $ 1 \times n$ board?  What are the initial conditions?
\vfill


}


\frame
{
  \frametitle{We need some recurrence relations}


{\bf  Problem 1}: How many flags with $n$ horizontal stripes can you make with $3$ colors so that adjacent stripes are different colors and the top and bottom stripes are different?
\vspace{.25cm}

{\bf Problem 2}: How many different tiling are there of a $1\times n$ board with $1 \times 1$ white tiles, $2 \times 1$ white tiles, and $2 \times 1$ black tiles?
\vspace{.25cm}

{\bf Problem 3}: How many different tiling are there of a $2\times n$ board with $1 \times 1$  tiles, and right triominos?

}

\frame
{
  \frametitle{We need some more recurrence relations}

%All Jacobsthal sequence

{\bf Problem 1}: How many different tiling are there of a $1\times n$ board with $1 \times 1$ white tiles, $2 \times 1$ white tiles, and $2 \times 1$ black tiles?
\vspace{.25cm}

{\bf Problem 2}: How many different tiling are there of a $2\times n$ board with $2 \times 1$  tiles, and $2 \times 2$  tiles?

{\bf Problem 3}: How many different tiling are there of a $3\times n$ board with $1 \times 1$  tiles, and $2 \times 2$  tiles?
\vspace{.25cm}

{\bf Problem 4}: How many ternary strings of length $n$ are there with no  two consecutive letters are non-zero?
\vspace{.25cm}

}


\frame
{
  \frametitle{Flags}

How many different flags $f_n$ with $n$ horizontal stripes can you make with $3$ colors so that adjacent stripes are different colors?  What is $f_1$, $f_2$ and $f_3$? Find a recurrence relation using the addition principle.  Find a closed form solution with the multiplication principle.

  \begin{figure}[htp]
  \begin{center}
    \subfigure{\includegraphics[scale=.45]{flags}}
    \end{center}
\end{figure}

}

\frame
{
  \frametitle{Staircases and Down-Right Paths}



  \begin{figure}[htp]
  \begin{center}
    \subfigure{\includegraphics[scale=.38]{Catalan_stairstep_5}}
    \subfigure{\includegraphics[scale=.135]{dyck_path_5}}
    \end{center}
\end{figure}

}

\frame
{
  \frametitle{$n=6$ Case}
\begin{figure}[htp]
    \subfigure{\includegraphics[scale=.75]{66board}}
\end{figure}

}

\end{comment}

\end{document}

